\documentclass[11pt]{article}
\usepackage[margin=25mm]{geometry}
\usepackage{parskip}
\pagestyle{empty}
\begin{document}

\noindent Ku Kang, PhD\\
CBRN Defense Research Institute\\
Republic of Korea\\
bisu9082@mnd.go.kr\\
ORCID 0000-0001-8739-218X

\vspace{6mm}
\noindent The Editors\\
\textit{Water Research}

\vspace{6mm}
\noindent Dear Editors,

We submit \textbf{``Testing effort sets the water-quality criterion: a
support-dependent artefact that regulatory screening narrows but does not
remove''} for consideration as a Research Paper.

Every major framework for deriving an aquatic water-quality criterion takes an
extreme of the available toxicity data: the lowest no-effect concentration under
the European guidance, the four lowest genus means under the United States
approach, the fifth percentile of a species sensitivity distribution under the
Canadian and Australian protocols. An extreme is an order statistic. Its
expected value moves with how many observations exist, and testing effort across
chemicals is anything but uniform. This paper asks what that does to the
criteria themselves and to the water bodies they are applied to.

The core of the work is a water-side result. We derive criteria for PFAS from a
single curated ECOTOX extract under seven conventions and join them to 20,682
measured concentrations from 1180 United States stations. Because the
conventions do not all reach the same substances, we report the comparison on
the six substances and 10,362 measurements every convention can assess: there
the fraction of samples exceeding the criterion runs from 98.8\,\% to zero on
the derivation choice alone. Perfluorooctanoic acid and perfluoroheptanoic
acid, adjacent homologues measured at the same 810 stations at comparable
concentrations, receive fully screened criteria 523-fold apart, behind which
stand 313 and 6 apical chronic records.

We then tested whether careful practice removes the problem, because that is
the obvious objection. We reproduced the screen regulatory agencies apply ---
apical endpoints, chronic exposure durations, active ingredient, species mean
chronic values, a log-normal species sensitivity distribution. It works in the
sense that matters for credibility: our fifth percentile for PFOS is
4738\,ng\,L$^{-1}$ against Canada's independently derived guideline of
6800\,ng\,L$^{-1}$, the closest of five available comparisons. And it does
real work on the disagreement between derivation routes, bringing the two
agencies actually use from 78.6\,\% against 4.6\,\% of samples to 14.8\,\%
against 4.6\,\%. What it does not touch is the dependence they share
($\rho = -0.590$ and $-0.681$). Screening reconciles the routes without
removing the bias in either.

A second finding came out of asking why. For the minimum-anchored conventions
the mechanism is ordinary order statistics. For the species sensitivity
distribution it is not: an HC5 is a parametric functional and should carry no
asymptotic dependence on sample size, and the dependence we observe enters
because the fitted spread widens as a substance is tested across more taxa and
conditions. That channel supplies 63\,\% of the movement. The distinction
matters because it sends any correction to a different place.

To establish that this is not a property of PFAS or of our substance selection,
we put the same question to all 4757 chemicals in ECOTOX with freshwater acute
data. Equalising sampling depth on species removes 50 to 61\,\% of the
minimum's correlation with record count, and controlling in addition for
within-chemical spread leaves a partial correlation of $-0.004$. We had
expected a residual association between hazard and attention to survive this
and it does not: the two channels account for essentially all of it.

The work bears directly on recent work in this journal on the conservatism of
the assessment-factor method (Ding et al., \textit{Water Research} 290,
125054). That programme treats the conservatism as a bias of roughly fixed size
to be corrected by better-calibrated factors. Our results indicate it is not
fixed: it varies with the evidential support behind each substance, so a
uniform correction will over-correct the data-poor compounds and under-correct
the data-rich ones. The analysis is built on public monitoring and public
toxicity data, and its consequence is expressed in stations rather than in
model statistics.

All analysis code, derived tables and the characteristic-name to structure
mapping with verification formulae are available in a public repository, and the
pipeline carries three assertion gates that reproduce previously deposited
tables before any new result is computed.

This manuscript is original, has not been published elsewhere, and is not under
consideration by another journal. All authors have approved the submission and
declare no competing financial or personal interests. Use of AI tools in
preparing the manuscript is declared in the manuscript itself, in the section
required by the journal's policy.

\vspace{4mm}
\noindent Yours sincerely,

\vspace{6mm}
\noindent Ku Kang\\
on behalf of Jeongyun Kim and Jin Yoo

\end{document}
